% Companion long-record migration for Erdős #269 round 6.
% Destination label: supporting extract of original sec:lcm (compressed identity retained in the short note)
% Not frozen R. Not a blind \input. Labels may collide with the short note;
% assign destination labels as specified in R6LongRecordMigrations.md.
% Extracted from Type B edits/relocation_extracts.tex.

\section{The running least common multiple as a product of pure powers}\label{sec:lcm}

\begin{proposition}[the running least common multiple]\label{res:lcm}
Let $p,q,r$ be pairwise distinct primes and $x\ge1$.  Then
$\Lprefix(x)=\hgt(x)$.
\end{proposition}

\begin{proof}
For divisibility in one direction, every smooth $n\le x$ has exponents bounded
by the corresponding integer logarithms, so $n\mid\hgt(x)$, and hence
$\Lprefix(x)\mid\hgt(x)$.  For the other, the three pure powers
$p^{\lfloor\log_p x\rfloor}$, $q^{\lfloor\log_q x\rfloor}$ and
$r^{\lfloor\log_r x\rfloor}$ are themselves smooth numbers not exceeding $x$,
so each divides $\Lprefix(x)$.  Distinct primes have coprime powers, so their
product divides $\Lprefix(x)$ as well.  The two divisibilities give equality.
\end{proof}

\begin{proposition}\label{res:cube}
For every $x\ge1$, $\hgt(x)\le x^{3}$.
\end{proposition}

\begin{proof}
Each of the three factors is a power of its base not exceeding $x$.
\end{proof}

The opposite estimate follows from
$x/b<b^{\lfloor\log_b x\rfloor}$ for each $b\in\{p,q,r\}$:
\[
 \frac{x^3}{pqr}<\hgt(x)\le x^3.
\]
In particular the reciprocal kernel is summable over the exponent lattice.
